% Copyright (C) 2014-2025 by Thomas Auzinger <thomas@auzinger.name>

% Inside the square brackets (e.g., [draft,onside]), different options can be chosen:
% 'draft': Obtain debug information on the build.
% 'final': Remove all debug information in the final build.
% 'onside': Create document for onsided printing (without empty pages added by LaTeX).
% 'twoside': Create document for twosided printing, which adds empty pages to start chapters on odd pages.
\documentclass[draft,twoside]{vutinfth}

% Define convenience functions to use the author name and the thesis title in the PDF document properties.
\newcommand{\authorname}{Sami Pasquali} % The author name without titles.
\newcommand{\thesistitle}{Design Games} % The title of the thesis. The English version should be used, if it exists.

% Create the XMP metadata file for the creation of PDF/A compatible documents.
\begin{filecontents*}[overwrite]{\jobname.xmpdata}
\Author{\authorname}                                    % The author's name in the document properties.
\Title{\thesistitle}                                    % The document's title in the document properties.
\Language{de-AT}                                        % The document's language in the document properties. Select 'en-US', 'en-GB', or 'de-AT'.
\Keywords{a\sep list\sep of\sep keywords}               % The document's keywords in the document properties (separated by '\sep ').
\Publisher{TU Wien}                                     % The document's publisher in the document properties.
\Subject{Thesis}                                        % The document's subject in the document properties.
\end{filecontents*}

% Load packages to allow in- and output of non-ASCII characters.
\usepackage{lmodern}        % Use an extension of the original Computer Modern font to minimize the use of bitmapped letters.
\usepackage[T1]{fontenc}    % Determines font encoding of the output. Font packages have to be included before this line.
\usepackage[utf8]{inputenc} % Determines encoding of the input. All input files have to use UTF8 encoding.

% Extended LaTeX functionality is enables by including packages with \usepackage{...}.
\usepackage{amsmath}    % Extended typesetting of mathematical expression.
\usepackage{amssymb}    % Provides a multitude of mathematical symbols.
\usepackage{mathtools}  % Further extensions of mathematical typesetting.
\usepackage{microtype}  % Small-scale typographic enhancements.
\usepackage[inline]{enumitem} % User control over the layout of lists (itemize, enumerate, description).
\usepackage{multirow}   % Allows table elements to span several rows.
\usepackage{booktabs}   % Improves the typesetting of tables.
\usepackage{subcaption} % Allows the use of subfigures and enables their referencing.
\usepackage[ruled,linesnumbered,algochapter]{algorithm2e} % Enables the writing of pseudo code.
\usepackage[dvipsnames,table]{xcolor} % Allows the definition and use of colors. This package has to be included before tikz.
\usepackage{nag}        % Issues warnings when best practices in writing LaTeX documents are violated.
\usepackage{todonotes}  % Provides tooltip-like todo notes.
\usepackage{morewrites} % Increases the number of external files that can be used.
\usepackage[a-2b,mathxmp]{pdfx}      % Enables PDF/A compliance. Loads the package hyperref and has to be included second to last.
\usepackage[acronym,toc]{glossaries} % Enables the generation of glossaries and lists of acronyms. This package has to be included last.

% Set PDF document properties
\hypersetup{
    pdfpagelayout   = TwoPageRight,           % How the document is shown in PDF viewers (optional).
    linkbordercolor = {Melon},                % The color of the borders of boxes around hyperlinks (optional).
}

\setpnumwidth{2.5em}        % Avoid overfull hboxes in the table of contents (see memoir manual).
\setsecnumdepth{subsection} % Enumerate subsections.

\nonzeroparskip             % Create space between paragraphs (optional).
\setlength{\parindent}{0pt} % Remove paragraph indentation (optional).

\makeindex      % Use an optional index.
\makeglossaries % Use an optional glossary.
%\glstocfalse   % Remove the glossaries from the table of contents.

% Set persons with 4 arguments:
%  {title before name}{name}{title after name}{gender}
%  where both titles are optional (i.e. can be given as empty brackets {}).
\setauthor{}{\authorname}{}{}
\setadvisor{Pretitle}{Hilda Tellioğlu}{Posttitle}{female}

% Required data.
\setregnumber{11809912}
\setdate{30}{09}{2026} % Set date with 3 arguments: {day}{month}{year}.
\settitle{\thesistitle}{Design Games} % Sets English and German version of the title (both can be English or German). If your title contains commas, enclose it with additional curvy brackets (i.e., {{your title}}) or define it as a macro as done with \thesistitle.
\setsubtitle{}{} % Sets English and German version of the subtitle (both can be English or German).

% Select the thesis type: bachelor / master / doctor.
% Bachelor:
\setthesis{bachelor}
%
% Master:
%\setthesis{master}
%\setmasterdegree{dipl.} % dipl. / rer.nat. / rer.soc.oec. / master
%
% Doctor:
%\setthesis{doctor}
%\setdoctordegree{rer.soc.oec.}% rer.nat. / techn. / rer.soc.oec.

% For bachelor and master:
\setcurriculum{Software Engineering}{Software Engineering} % Sets the English and German name of the curriculum.

% Optional reviewer data:
%\setfirstreviewerdata{Affiliation, Country}
%\setsecondreviewerdata{Affiliation, Country}


\begin{document}

\frontmatter % Switches to roman numbering.
% The structure of the thesis has to conform to the guidelines at
%  https://informatics.tuwien.ac.at/study-services

\addtitlepage{naustrian} % German title page.
\addtitlepage{english} % English title page.
\addstatementpage{naustrian}  % Select 'english' for the English statement text.

\begin{danksagung*}
\todo{Ihr Text hier.}
\end{danksagung*}

\begin{acknowledgements*}
\todo{Enter your text here.}
\end{acknowledgements*}

\begin{kurzfassung}
\todo{Ihr Text hier.}
\end{kurzfassung}

\begin{abstract}
\todo{Enter your text here.}
\end{abstract}

% Select the language of the thesis, e.g., english or naustrian.
\selectlanguage{english}

% Add a table of contents (toc).
\tableofcontents % Starred version, i.e., \tableofcontents*, removes the self-entry.

% Switch to arabic numbering and start the enumeration of chapters in the table of content.
\mainmatter

\chapter{Design Games}
\cite{brandtDesigningExploratoryDesign2006} describes Design games as a way of organizing participation in design projects. \cite{triantafyllakosDesigningEducationalSoftware2011} notes that they (Design Games) provide a starting point for facilitating user participation and contribution. 
\subsection{History}
\subsection{Connection to Games}
\subsection{Categories and Classification of Design Games}

\subsection{Mechanics}
\subsection{Benefits}
\subsection{Problems and Issues}
\chapter{Examples of Design Game and their use in different domains}
\chapter{Design Games in Design Thinking}
\subsection{Opportunities for Design Game use}

\backmatter

% Declare the use of AI tools as mentioned in the statement of originality.
% Use either the English aitools or the German kitools.
\begin{aitools}
None
\end{aitools}

\begin{kitools}
Keine
\end{kitools}

% Use an optional list of figures.
\listoffigures % Starred version, i.e., \listoffigures*, removes the toc entry.

% Use an optional list of tables.
\cleardoublepage % Start list of tables on the next empty right hand page.
\listoftables % Starred version, i.e., \listoftables*, removes the toc entry.

% Use an optional list of algorithms.
\listofalgorithms
\addcontentsline{toc}{chapter}{List of Algorithms}

% Add an index.
\printindex % Rerun build-thesis or build-all when removing the index through commenting.

% Add a glossary.
\printglossaries % Rerun build-thesis or build-all when removing it through commenting.

% Add a bibliography.
\bibliographystyle{alpha}
\bibliography{references}

\end{document}